%% LunarAtlas: Draft with LTTB Approach (Revised Manuscript)
%% Elsevier two-column style (cas-dc class)
%% Authors: Lovekesh Anand, Dua Saeed
%% April 2026

\documentclass[a4paper,twocolumn,10pt]{article}
\usepackage[margin=1.9cm,top=2.6cm,bottom=2.6cm,columnsep=0.6cm]{geometry}
\usepackage{authblk}

% -----------------------------------------------------------------------
%  Core packages
% -----------------------------------------------------------------------
\usepackage{amsmath,amssymb}
\usepackage{siunitx}
\usepackage{graphicx}
\usepackage{booktabs}
\usepackage{tabularx}
\usepackage{multirow}
\usepackage{hyperref}
\usepackage{url}
\usepackage{enumitem}
\usepackage{subcaption}
\usepackage{float}
\usepackage{xcolor}
\usepackage{microtype}
\usepackage{lineno}   % line numbers for review
\usepackage{stfloats} % allows figure* at bottom of page
\usepackage[font=small,labelfont=bf,justification=centering]{caption}

% -----------------------------------------------------------------------
%  SI units configuration
% -----------------------------------------------------------------------
\sisetup{
  detect-all,
  group-separator={,},
  separate-uncertainty=true
}

% -----------------------------------------------------------------------
%  Hyperref colours (Elsevier-like)
% -----------------------------------------------------------------------
\hypersetup{
  colorlinks=true,
  linkcolor=blue!60!black,
  citecolor=blue!60!black,
  urlcolor=blue!60!black
}

% -----------------------------------------------------------------------
%  Elsevier-like title/abstract styling
% -----------------------------------------------------------------------
\usepackage{titlesec}
\usepackage{fancyhdr}
\usepackage{abstract}
\usepackage{etoolbox}

% Section formatting — Elsevier style
\titleformat{\section}{\normalfont\bfseries\fontsize{10}{12}\selectfont}{\thesection.}{0.5em}{}
\titleformat{\subsection}{\normalfont\bfseries\fontsize{9.5}{11}\selectfont}{\thesubsection}{0.5em}{}
\titleformat{\subsubsection}{\normalfont\itshape\fontsize{9}{11}\selectfont}{\thesubsubsection}{0.5em}{}
\titlespacing*{\section}{0pt}{10pt plus 2pt}{4pt plus 1pt}
\titlespacing*{\subsection}{0pt}{8pt plus 1pt}{3pt plus 1pt}
\titlespacing*{\subsubsection}{0pt}{6pt}{2pt}

% Header/footer
\pagestyle{fancy}
\fancyhf{}
\fancyhead[L]{\small\textit{L.~Anand, D.~Saeed / LunarAtlas: Chandrayaan-3 LIBS Infrastructure}}
\fancyhead[R]{\small\textit{April 2026}}
\fancyfoot[C]{\small\thepage}
\renewcommand{\headrulewidth}{0.3pt}

% Abstract styling
\renewcommand{\abstractnamefont}{\normalfont\bfseries\small}
\renewcommand{\abstracttextfont}{\normalfont\small}
\setlength{\absleftindent}{0pt}
\setlength{\absrightindent}{0pt}

% Paragraph spacing
\setlength{\parskip}{3pt}
\setlength{\parindent}{12pt}

% =========================================================================
%  Document begins
% =========================================================================
\begin{document}

\twocolumn[{%
\begin{@twocolumnfalse}

{\centering
{\LARGE\bfseries LunarAtlas: From PDS4 Archives to\\[4pt]
Analysis-Ready Spectra --- A Reproducible\\[4pt]
Infrastructure for Chandrayaan-3 LIBS Data}\par
\vspace{8pt}
{\normalsize\bfseries Lovekesh Anand\textsuperscript{a},\quad Dua Saeed\textsuperscript{a}}\\[3pt]
{\small\textit{\textsuperscript{a}Independent Researchers, India}}
\par
}

\vspace{10pt}
\begin{abstract}
\noindent
We present \textbf{LunarAtlas}, a reproducible data-processing
infrastructure that transforms publicly available Chandrayaan-3
Laser-Induced Breakdown Spectroscopy (LIBS) Level-1 (L1) products
into cleaned, machine-accessible, long-form spectral records suitable
for quantitative lunar science.
Starting from calibrated L1 tables released through ISRO's
PDS4-compliant archive, LunarAtlas implements a transparent Python
pipeline that: (i)~parses XML metadata; (ii)~reshapes wide-format
Chandrayaan-3 L1 tables---in which 2\,049 wavelength channels appear
as column headers spanning \SIrange{164.35}{878.26}{\nano\metre}---into
normalised per-wavelength long-form records; (iii)~identifies and
correctly pairs plasma and background measurements using mission flag
columns (\texttt{Laser\_Fire\_Status}, \texttt{Force\_Reset\_Status});
and (iv)~performs physically motivated background subtraction,
$I_{\text{clean}}(\lambda) = \max(0,\,I_{\text{plasma}}(\lambda)
- I_{\text{background}}(\lambda))$, with non-physical negatives
clamped to zero.
Each processed file receives an MD5 checksum and an algorithm version
record.
Applied to \texttt{ch3\_lib\_002\_20230825T104221\_00\_l1.csv}, the
pipeline yields four cleaned measurements (\textbf{Measurement~IDs~1--4}),
each with 2\,094 wavelength channels and 19 metadata columns.
A baseline suppression factor of $\sim$7.1 and a high-frequency
variance increase of $\sim$3.9 are observed, validating the physical
correctness of the pipeline.
LunarAtlas further provides a PDS4-aware PostgreSQL schema and implements
an adaptive \textbf{Largest Triangle Three Buckets (LTTB)} downsampling
strategy with mathematically defined zoom-saturation criteria that
guarantees preservation of all emission-line peaks while achieving
sub-\SI{500}{\milli\second} response times.  The LTTB algorithm operates
over four distinct zoom levels ($z=0$ to $z=3$), progressively revealing
spectral detail from full-scan overview to raw-resolution feature
inspection, with automatic fallback to raw data when zoom exceeds
instrumental resolution limits.
By establishing a ``single source of scientific truth,'' the system
offers a replicable methodology for Chandrayaan-4, Artemis-era VIPER,
and Commercial Lunar Payload Services (CLPS) LIBS instruments.
\end{abstract}
\vspace{4pt}
\noindent\textbf{Keywords:} Laser-Induced Breakdown Spectroscopy;
Chandrayaan-3; Lunar Regolith; PDS4; Data Cleaning; Spectral
Reproducibility; Background Subtraction; LTTB Downsampling;
Measurement~ID.

\vspace{10pt}\hrule\vspace{12pt}
\end{@twocolumnfalse}
}]

% \linenumbers  % disabled for submission version

% =========================================================================
\section{Introduction}
\label{sec:intro}
% =========================================================================

\subsection{The Promise of Lunar In-Situ Spectroscopy}

The elemental composition of the lunar surface encodes four billion
years of planetary evolution---from the primordial magma ocean, through
impact gardening, to ongoing solar-wind implantation of hydrogen and
helium.  Accessing this record has long depended on remote sensing and
the painstaking laboratory analysis of returned samples.  Laser-Induced
Breakdown Spectroscopy (LIBS) offers a fundamentally different path: a
pulsed laser focused on a surface target ablates a few nanograms to
micrograms of material, generating a transient plasma whose
characteristic optical emission encodes the elemental abundances of the
ablated material with no sample preparation and no contact requirement
\cite{Cremers2006,Wiens2012}.

The technique was first demonstrated in a planetary setting aboard
NASA's Mars Science Laboratory (MSL) rover \emph{Curiosity} in 2012,
where the ChemCam instrument acquired more than 930\,000 single-shot
spectra from 3\,500 unique rock and regolith targets in Gale
Crater~\cite{Wiens2012,Maurice2012,Gasnault2023}.  The global
community subsequently adopted LIBS for the SuperCam instrument on
NASA's \emph{Perseverance} rover and the MarSCoDe instrument on
China's Zhurong rover, cementing LIBS as the \emph{de facto} in-situ
geochemical sensor for planetary surfaces~\cite{Wiens2020,Xu2021}.

Against this backdrop, the successful soft landing of India's
Chandrayaan-3 mission near the lunar south pole on 23~August~2023
marked a milestone of particular significance: the \emph{first
deployment of LIBS on the Moon}~\cite{ISROCh3}.  The Pragyan rover
carried a miniaturised LIBS instrument alongside an Alpha Particle
X-ray Spectrometer (APXS)~\cite{Laxmiprasad2020}.  Operating in the
near-vacuum lunar environment introduces plasma dynamics fundamentally
different from the $\sim$7~mbar Martian atmosphere where ChemCam was
optimised~\cite{Saeidfirozeh2024}, underscoring the scientific
importance of correctly handling and interpreting the resulting spectral
data.

\subsection{The Data Accessibility Gap}

ISRO released the Chandrayaan-3 LIBS data as calibrated Level-1 (L1)
products conforming to Planetary Data System version~4 (PDS4)
standards~\cite{Hughes2018}.  PDS4 is a rigorous archival format: each
data product consists of a self-describing XML label paired with a
tabular data file, ensuring long-term preservation and provenance.

However, the Chandrayaan-3 L1 tables present a \emph{structural
challenge} that is not immediately obvious from the PDS4 documentation:
the intensity data are encoded in a \emph{wide format} in which the
2\,049 wavelength channels (spanning \SIrange{164.35}{878.26}{\nano
\metre}) appear as column \emph{headers}, with their numeric wavelength
values used directly as column names (e.g., \texttt{164.35},
\texttt{164.74}, \ldots, \texttt{878.26}).  This wide-format structure
is not directly amenable to relational database storage, spectral queries,
pairing of plasma and background shots, or downstream peak detection.

Furthermore, interactive visualisation of these spectra at multiple zoom
levels requires intelligent downsampling that preserves narrow emission
lines---a requirement not met by standard time-series algorithms that
prioritise envelope preservation over peak retention.

\subsection{Research Question and Contributions}

The central question addressed by this paper is:

\begin{quote}
\emph{How can Chandrayaan-3 LIBS Level-1 data be cleaned, normalised,
downsampled, and structured into a reproducible, queryable, and
machine-accessible resource without compromising spectral fidelity
or peak integrity?}
\end{quote}

This paper makes the following contributions, ordered by scientific
importance:

\begin{enumerate}[label=\textbf{(C\arabic*)},leftmargin=*,itemsep=3pt]
  \item \textbf{Level-1 cleaning pipeline.}  A transparent,
    step-by-step Python pipeline that ingests Chandrayaan-3 LIBS L1
    tables, correctly identifies plasma versus background measurements
    using mission flag columns, reshapes wide-format data into
    normalised long-form spectra, and performs physically motivated
    background subtraction.

  \item \textbf{Measurement~ID semantics.}  A formal definition and
    implementation of the \emph{Measurement~ID}---the key identifier
    that links each cleaned spectral record to its originating plasma
    shot, instrument configuration, and acquisition timestamp.

  \item \textbf{PDS4-aware data model.}  A normalised PostgreSQL
    schema that stores cleaned spectral samples together with the
    complete mission hierarchy, with MD5 file checksums and
    algorithm-version metadata for every ingested file.

  \item \textbf{LTTB-based adaptive downsampling.}  A multi-zoom
    \textbf{Largest Triangle Three Buckets (LTTB)} algorithm adapted
    for fixed-channel LIBS spectra, with dual zoom-saturation criteria
    and guaranteed peak preservation across all zoom levels.

  \item \textbf{Empirical cleaning validation.}  Quantitative
    characterisation of the effect of background subtraction on a real
    Chandrayaan-3 L1 file, including baseline suppression factors and
    high-frequency variance behaviour, reported for the first time for
    this specific mission product.
\end{enumerate}

\subsection{Paper Organisation}

Section~\ref{sec:related} reviews related work.
Section~\ref{sec:data} describes the Chandrayaan-3 dataset and PDS4
context in detail.  Section~\ref{sec:architecture} presents the
LunarAtlas architecture and data model.
Section~\ref{sec:processing} details the Level-1 cleaning pipeline.
Section~\ref{sec:lttb} introduces the LTTB downsampling formulation.
Section~\ref{sec:api} describes the API and visualisation workflow.
Section~\ref{sec:results} presents cleaning results and LTTB validation.
Section~\ref{sec:discussion} discusses implications and limitations.
Section~\ref{sec:conclusion} concludes.


% =========================================================================
\section{Related Work}
\label{sec:related}
% =========================================================================

\subsection{Planetary LIBS Instruments and Mission Heritage}

The deployment of LIBS on planetary surfaces has a rich history
beginning with ChemCam on MSL \emph{Curiosity}~\cite{Wiens2012,
Maurice2012}.  ChemCam's three Czerny-Turner spectrometers cover
240--342~nm (UV), 382--469~nm (violet), and 474--907~nm (VNIR), and
its Nd:KGW laser at 1067~nm has enabled remote analysis of targets up
to 7~m distant.  SuperCam on \emph{Perseverance} extends ChemCam's
capability by integrating LIBS with Raman spectroscopy~\cite{Wiens2020}.
China's MarSCoDe instrument on Zhurong combines LIBS and SWIR
spectroscopy across 240--2400~nm~\cite{Xu2021,Zhao2023}.

For lunar LIBS, earlier feasibility work demonstrated acceptable
signal-to-noise in simulated vacuum conditions~\cite{Knight2000,
Lasue2012}.  The Chandrayaan-3 LIBS instrument represents a miniaturised
system weighing approximately 1.1~kg with power consumption below
1.2~W~\cite{Laxmiprasad2020}.

\subsection{Downsampling for Interactive Spectral Visualisation}

General-purpose time-series downsampling algorithms---including mean
aggregation, Douglas-Peucker curve simplification, and the Largest
Triangle Three Buckets (LTTB) method~\cite{Steinarsson2013}---are
widely deployed in web visualisation frameworks.  LTTB was originally
developed for financial and sensor time-series data and operates by
dividing the data into buckets and selecting the point within each
bucket that forms the largest triangle with neighbouring points,
thereby preserving visual shape while dramatically reducing point count.

However, standard LTTB implementations do not guarantee preservation
of narrow emission lines, which may span only one to three wavelength
channels.  For LIBS spectroscopy, where a Ca~II doublet at 393--397~nm
or an Mg~II line at 280~nm can be the primary diagnostic signature,
we require an \emph{enhanced LTTB} strategy that explicitly unions
LTTB-selected points with all local maxima, ensuring 100\% peak
retention across all zoom levels.


% =========================================================================
\section{Chandrayaan-3 LIBS Dataset and PDS4 Context}
\label{sec:data}
% =========================================================================

\subsection{Mission Overview}

Chandrayaan-3 was launched on 14~July~2023 by the Indian Space Research
Organisation (ISRO) and achieved a successful soft landing near the
lunar south pole at 69.37$^\circ$S, 32.35$^\circ$E on 23~August~2023,
at a site subsequently named \emph{Shiv Shakti Point}~\cite{ISROCh3}.
The lander, \emph{Vikram}, deployed the \emph{Pragyan} rover, which
carried two in-situ science payloads: the LIBS instrument and the APXS.

\subsection{The Pragyan LIBS Instrument}

The LIBS instrument aboard Pragyan employs a pulsed Nd:YAG laser
focused on the lunar surface at close range and collects plasma
emission with a spectrometer spanning \SIrange{164.35}{878.26}{\nano
\metre}.  The operational measurement sequence consists of alternating
\emph{background} acquisitions (laser not fired) and \emph{plasma}
acquisitions (laser fired), allowing background subtraction to isolate
true LIBS emission.

\subsection{Level-1 Data Format: Wide-Table Structure}

The Chandrayaan-3 LIBS L1 products are distributed in a wide-format
tabular structure.  Table~\ref{tab:l1_schema} summarises the column
organisation.

\begin{table}[htbp]
\centering
\caption{Column organisation of Chandrayaan-3 LIBS Level-1 CSV files.}
\label{tab:l1_schema}
{\small
\begin{tabularx}{\columnwidth}{@{}lX@{}}
\toprule
\textbf{Group} & \textbf{Contents} \\
\midrule
Metadata~(6) & Time, Measurement Count, Operation Mode,
  Measurement Type, Force Reset Status, Laser Fire Status \\
\addlinespace
Wavelengths~(2\,049) & Column headers are wavelength values
  in nm; cells are intensity counts \\
\addlinespace
Housekeeping~(8) & Delay Time, Integration Time,
  Number of Pulses, X-Factor, Laser Energy,
  Pump Current, PRR, On/Off Status \\
\bottomrule
\end{tabularx}
}
\end{table}

Each row represents a single measurement acquisition.  The file
\texttt{ch3\_lib\_002\_20230825T104221\_00\_l1.csv} contains eight
rows: four background and four plasma acquisitions.

\subsection{Plasma and Background Flag Semantics}

Background and plasma measurements are distinguished by two binary
flag columns:
\begin{itemize}[nosep,leftmargin=*]
  \item \texttt{Force\_Reset\_Status}: value 1 indicates background;
    value 0 indicates normal operation.
  \item \texttt{Laser\_Fire\_Status}: value 0 for background; value 1
    for plasma.
\end{itemize}

A valid plasma measurement satisfies both
\texttt{Force\_Reset\_Status}~$=0$ \emph{and}
\texttt{Laser\_Fire\_Status}~$=1$.


% =========================================================================
\section{LunarAtlas System Architecture}
\label{sec:architecture}
% =========================================================================

\subsection{Design Principles}

LunarAtlas is structured around three architectural principles:

\begin{enumerate}[nosep,leftmargin=*]
  \item \textbf{Single source of scientific truth.}  All domain logic
    resides exclusively in the backend.  The visualisation layer
    performs no scientific computation.

  \item \textbf{Reproducibility by design.}  Every processing step
    records its algorithm version, timestamp, and MD5 checksum.

  \item \textbf{Separation of concerns.}  Data persistence, scientific
    logic, and presentation are independently testable.
\end{enumerate}

\subsection{Three-Tier Architecture}

The system comprises three tiers:

\textbf{Data layer (PostgreSQL).}  The relational database stores the
complete mission hierarchy with MD5 checksums and algorithm versions.

\textbf{Logic layer (Python / FastAPI).}  The backend implements all
scientifically meaningful operations: XML label parsing, wide-to-long
reshaping, background subtraction, LTTB downsampling, peak detection,
and NIST line matching.

\textbf{Presentation layer (React~+ amCharts).}  The frontend manages
viewport state and chart rendering, issuing HTTP requests for each
viewport change.


% =========================================================================
\section{Data Processing and Cleaning Pipeline}
\label{sec:processing}
% =========================================================================

\subsection{Pipeline Overview}

The LunarAtlas cleaning pipeline is implemented in Python and consists
of two primary modules: \texttt{batch\_process\_libs.py} for data
ingestion and cleaning, and \texttt{plot\_libs\_spectra.py} for
diagnostic visualisation.

\subsection{Step 1: XML Label Parsing and Metadata Extraction}

The pipeline parses the PDS4 XML label using Python's
\texttt{xml.etree.ElementTree} module, extracting: PDS4 logical
identifier, observation timestamps, source file name and size, column
descriptions, and instrument context.

\subsection{Step 2: File Integrity Verification (MD5)}

Before processing, the pipeline computes an MD5 checksum of the source
L1 file, enabling corruption detection, traceability, and deduplication.

\subsection{Step 3: Wide-to-Long Reshaping}

The wide DataFrame is reshaped using a \texttt{melt} operation: each
row produces 2\,049 long-form rows, one per wavelength channel.  The
resulting schema comprises measurement keys, spectral data
(wavelength, intensity), and housekeeping parameters.

\subsection{Step 4: Plasma and Background Identification}

Using the flag columns, the pipeline separates records into background
and plasma sub-DataFrames:

\begin{align}
\text{Background:}\quad & F = 1 \;\wedge\; L = 0 \\
\text{Plasma:}\quad & F = 0 \;\wedge\; L = 1
\end{align}

where $F$ denotes \texttt{Force\_Reset\_Status} and $L$ denotes
\texttt{Laser\_Fire\_Status}.

\subsection{Step 5: Background Subtraction and Negative-Value Clamping}

For each background--plasma pair, the cleaned spectrum is:

\begin{equation}
\label{eq:subtraction}
I_{\text{clean}}(\lambda) =
  \max\!\bigl(0,\;
  I_{\text{plasma}}(\lambda) - I_{\text{background}}(\lambda)
  \bigr).
\end{equation}

The $\max(0, \cdot)$ operator clamps negative values to zero, as
negative photon counts are physically meaningless.

\subsection{The Measurement~ID: Definition and Significance}

The \textbf{Measurement~ID} is a sequential integer that uniquely
identifies a single cleaned spectral record within a given L1 file.
It groups the 2\,094 per-wavelength rows of a single acquisition,
serves as the foreign key linking spectral data to measurement metadata,
and enables per-measurement operations via simple \texttt{groupby}
operations.

For the archival file, four Measurement~IDs are produced
(Table~\ref{tab:measurement_ids}).

\begin{table}[htbp]
\centering
\caption{Measurement~ID characteristics for
  \texttt{ch3\_lib\_002\_20230825T104221\_00\_l1}.}
\label{tab:measurement_ids}
\begin{tabular}{@{}ccccc@{}}
\toprule
\textbf{ID} & \textbf{Time (UTC)} & \textbf{$E$ (V)} & \textbf{Peak~I} & \textbf{Peak~$\lambda$ (nm)} \\
\midrule
1 & 11:35:22 & 3327 & 321  & 748.0 \\
2 & 11:35:22 & 3280 & 261  & 283.1 \\
3 & 11:35:22 & 3194 & 684  & 282.3 \\
4 & 11:35:23 & 3202 & 637  & 282.7 \\
\bottomrule
\end{tabular}
\end{table}

Figure~\ref{fig:variability} shows the four cleaned measurements,
revealing dramatic shot-to-shot variability that would be masked by
bulk averaging.

\begin{figure}[htbp]
  \centering
  \includegraphics[width=\columnwidth]{fig5_variability.png}
  \caption{Cleaned spectra for Measurement~IDs~1--4, showing
    shot-to-shot variability. ID~3 exhibits high baseline across
    the full spectral range, while ID~2 shows near-complete
    background removal with emission concentrated near 283~nm.}
  \label{fig:variability}
\end{figure}


% =========================================================================
\section{LTTB-Based Adaptive Downsampling}
\label{sec:lttb}
% =========================================================================

\subsection{Problem Statement and LTTB Overview}

Let $S = \{(\lambda_i, I_i) : i = 1, \ldots, N\}$ be a cleaned
spectrum with $N = 2094$ channels.  For interactive visualisation,
we require a reduced representation $S'$ with $N' \ll N$ points that:
(i)~preserves all emission-line peaks; (ii)~maintains overall spectral
shape; and (iii)~adapts to user zoom level.

The Largest Triangle Three Buckets (LTTB) algorithm~\cite{Steinarsson2013}
addresses this by dividing the data into buckets and selecting the point
in each bucket that maximises the triangle area formed with the
previously selected point and the average of the next bucket.
Figure~\ref{fig:lttb_mechanism} illustrates the core mechanism.

\begin{figure}[htbp]
  \centering
  \includegraphics[width=\columnwidth]{fig1_lttb_mechanism.png}
  \caption{LTTB triangle area maximization. For each candidate bucket
    (cyan), the algorithm selects the point (red star) that forms the
    largest triangle with the previous selected point (blue) and the
    next bucket's centroid (black cross).}
  \label{fig:lttb_mechanism}
\end{figure}

\subsection{Multi-Zoom LTTB Formulation}

We implement a four-level zoom hierarchy ($z = 0, 1, 2, 3$):

\begin{itemize}[nosep,leftmargin=*]
  \item $z=0$: Full scan (coarsest downsampling, $\sim$150 points)
  \item $z=1$: Segment zoom (intermediate, $\sim$750 points)
  \item $z=2$: Feature zoom (fine detail, $\sim$42 points for narrow window)
  \item $z=3$: Raw resolution (no downsampling, all 2094 points)
\end{itemize}

The number of LTTB buckets at zoom level $z$ over wavelength span
$\Delta\lambda$ is:

\begin{equation}
N_{\text{buckets}}(z) = \text{BASE\_BUCKETS} \times 2^z,
\end{equation}

where $\text{BASE\_BUCKETS} \approx 150$ for $z=0$ full-scan view.

Figure~\ref{fig:adaptive_zoom} demonstrates the progressive zoom
behaviour across the three LTTB levels.

\begin{figure*}[!t]
  \centering
  \includegraphics[width=\textwidth]{fig3_adaptive_zoom.png}
  \caption{Adaptive zoom progression: (a)~Global view at $z=0$ shows
    full spectral range with LTTB downsampling; (b)~Regional zoom at
    $z=1$ reveals intermediate structure; (c)~Feature zoom at $z=2$
    preserves fine detail in narrow windows.}
  \label{fig:adaptive_zoom}
\end{figure*}

\subsection{Dual Zoom-Saturation Criteria}

Two independent constraints limit downsampling:

\begin{enumerate}[nosep,leftmargin=*]
  \item \textbf{Instrument resolution constraint}: Buckets cannot be
    narrower than the spectral resolution ($\sim$0.35~nm).
  \item \textbf{Data density constraint}: Cannot request more buckets
    than available raw samples in the current window.
\end{enumerate}

The zoom-saturation level $z_{\max}$ is:

\begin{equation}
z_{\max}(\Delta\lambda) =
  \min\!\left(
    \left\lfloor \log_2\!\left(
      \frac{\Delta\lambda}{\text{BASE\_BUCKETS} \times b_{\min}}
    \right)\right\rfloor,\;
    \left\lfloor \log_2\!\left(
      \frac{n_{\text{pts}}}{\text{BASE\_BUCKETS}}
    \right)\right\rfloor
  \right),
\end{equation}

where $b_{\min}$ is the minimum bucket width and $n_{\text{pts}}$ is
the number of raw points in the window.

For $z > z_{\max}$, the system automatically falls back to raw-data
mode, as shown in Figure~\ref{fig:zoom_saturation} panel~(d).

\begin{figure*}[!t]
  \centering
  \includegraphics[width=\textwidth]{_RN46CVM.png}
  \caption{Four-level zoom hierarchy demonstrating LTTB behaviour and
    saturation: (a)~$z=0$ full scan with 150~LTTB points;
    (b)~$z=1$ segment zoom with 750~points; (c)~$z=2$ feature zoom
    with 42~points over narrow window; (d)~$z=3$ raw resolution
    showing all available points when zoom exceeds $z_{\max}$.}
  \label{fig:zoom_saturation}
\end{figure*}

\subsection{Peak-Preservation Guarantee}

To ensure 100\% retention of emission lines, we compute the set of
local maxima:

\begin{equation}
P_{\text{peaks}} = \left\{ i \,\middle|\,
  I_i > I_{i-1} \;\wedge\; I_i > I_{i+1}
\right\},
\end{equation}

and take the union with LTTB-selected indices:

\begin{equation}
P_{\text{final}} = P_{\text{LTTB}} \cup P_{\text{peaks}}.
\end{equation}

This guarantees that no diagnostic emission line is omitted, regardless
of LTTB bucket boundaries.  Figure~\ref{fig:perfect_comparison} shows
that LTTB+peaks preserves structural detail at both global and zoomed
scales.

\begin{figure*}[!t]
  \centering
  \includegraphics[width=\textwidth]{fig2_perfect_comparison.png}
  \caption{LTTB structural preservation: (a)~Full spectrum comparison
    showing original (2094~pts, gray) versus LTTB (500~pts, blue)
    maintaining overall envelope; (b)~Detailed zoom demonstrating
    that LTTB+peaks (red path) captures all local maxima while
    following the original trace.}
  \label{fig:perfect_comparison}
\end{figure*}


% =========================================================================
\section{API Design and Visualisation Workflow}
\label{sec:api}
% =========================================================================

\subsection{Primary Spectral Endpoint}

The main API endpoint is:

\begin{verbatim}
GET /api/v1/spectral
\end{verbatim}

with query parameters: \texttt{observation\_id},
\texttt{min\_wavelength}, \texttt{max\_wavelength},
\texttt{zoom\_level}, and \texttt{include\_nist}.

\subsection{Progressive Loading Strategy}

The frontend follows a progressive loading protocol:

\begin{enumerate}[nosep,leftmargin=*]
  \item On initial load, request full-range spectrum at $z = 0$.
  \item On zoom/pan, extract new viewport $[\lambda_a, \lambda_b]$,
    compute $z_{\max}$, and request appropriate zoom level.
  \item At $z > z_{\max}$, backend returns raw samples.
\end{enumerate}

This ensures scientific fidelity at every zoom level with no
client-side interpolation.


% =========================================================================
\section{Results and Experimental Validation}
\label{sec:results}
% =========================================================================

\subsection{Cleaning Effectiveness: Baseline Suppression}

For a representative 20~nm window (560--580~nm):

\begin{itemize}[nosep,leftmargin=*]
  \item Raw: $\bar{I}_{\text{raw}} \approx 1005$~counts,
    $\sigma_{\text{raw}} \approx 47$~counts
  \item Cleaned: $\bar{I}_{\text{clean}} \approx 141$~counts,
    $\sigma_{\text{clean}} \approx 184$~counts
\end{itemize}

Baseline suppression factor: $\sim$7.1.
Variance increase: $\sim$3.9 (consistent with Poisson noise propagation).

\subsection{Inter-Measurement Variability}

Table~\ref{tab:measurement_stats} summarises cleaned-intensity
statistics across the four Measurement~IDs.

\begin{table}[htbp]
\centering
\caption{Cleaned intensity statistics per Measurement~ID.  The
  ``\% Neg.'' column reports the fraction of channels with negative
  values before clamping.}
\label{tab:measurement_stats}
{\small
\begin{tabular}{@{}cccccc@{}}
\toprule
\textbf{ID} & \textbf{Mean} & \textbf{Std} & \textbf{Max} & \textbf{\% Neg.} & \textbf{$E$~(V)} \\
\midrule
1 & 61.0  & 56.9  & 321  & 22.5 & 3327 \\
2 & 9.4   & 23.8  & 261  & 76.8 & 3280 \\
3 & 416.9 & 95.0  & 684  & 0.5  & 3194 \\
4 & 54.1  & 62.4  & 637  & 28.1 & 3202 \\
\bottomrule
\end{tabular}
}
\end{table}

The striking difference between ID~3 (mean~417~counts, 0.5\% negative)
and ID~2 (mean~9~counts, 76.8\% negative) highlights the importance
of per-shot Measurement~IDs.

\subsection{LTTB Downsampling Performance}

Across all zoom levels, LTTB+peaks achieves:

\begin{itemize}[nosep,leftmargin=*]
  \item 100\% peak retention (all local maxima preserved)
  \item Data reduction: 93\% at $z=0$ (150 vs 2094 points)
  \item Median API response time: $<$\SI{500}{\milli\second}
  \item Visual fidelity: indistinguishable from raw data at
    appropriate zoom levels
\end{itemize}


% =========================================================================
\section{Discussion}
\label{sec:discussion}
% =========================================================================

\subsection{The Reproducibility Gap in Planetary LIBS}

The cleaning and normalisation of Chandrayaan-3 LIBS data is not
merely a preprocessing convenience---it is a precondition for
reproducible science.  LunarAtlas addresses this by formalising the
wide-to-long reshaping, recording algorithm versions and MD5 checksums,
and providing an open, documented pipeline.

\subsection{LTTB Advantages for LIBS Spectroscopy}

The LTTB approach offers three key advantages over simple binning or
mean aggregation:

\begin{enumerate}[nosep,leftmargin=*]
  \item \textbf{Shape preservation}: Triangle-area maximization
    maintains visual structure at all scales.
  \item \textbf{Peak guarantee}: Explicit union with local maxima
    ensures 100\% retention of emission lines.
  \item \textbf{Adaptive zoom}: Dual saturation criteria prevent
    over-downsampling while enabling aggressive reduction where
    appropriate.
\end{enumerate}

\subsection{Generalisability to Other Missions}

The LunarAtlas architecture is deliberately general.  The core
operations apply to any mission delivering LIBS L1 products in
PDS4-compliant wide-format CSV.  This includes future Indian lunar
missions, asteroid LIBS instruments, and Mars missions releasing
raw L1 data.

\subsection{Limitations}

Current limitations include:

\begin{enumerate}[nosep,leftmargin=*]
  \item \textbf{Single-file validation.}  Quantitative metrics reported
    for one L1 file; full-archive validation needed.
  \item \textbf{Simple background pairing.}  Temporal-order pairing
    only; more sophisticated strategies may improve quality.
  \item \textbf{Absence of denoising.}  No SWT or Savitzky-Golay
    filtering; adding these would improve peak detection.
  \item \textbf{No user study.}  Practical utility of Measurement~ID
    and LTTB not evaluated with domain scientists.
\end{enumerate}


% =========================================================================
\section{Conclusion}
\label{sec:conclusion}
% =========================================================================

We have presented LunarAtlas, a reproducible data-processing
infrastructure that transforms Chandrayaan-3 LIBS Level-1 archives
into analysis-ready spectral records through transparent cleaning,
normalisation, and intelligent downsampling.

The core contributions are: (i)~a transparent Python pipeline for
background subtraction and wide-to-long reshaping; (ii)~the
Measurement~ID as a stable identifier for per-shot analysis;
(iii)~a PDS4-aware PostgreSQL schema with MD5 provenance; and
(iv)~an LTTB-based adaptive downsampling strategy with guaranteed
peak preservation across four zoom levels.

Applied to archival Chandrayaan-3 data, the system achieves baseline
suppression of $\sim$7.1$\times$, 100\% peak retention, and
sub-\SI{500}{\milli\second} API response times.  As the lunar
exploration community prepares for Chandrayaan-4 and Artemis missions,
LunarAtlas provides a working blueprint for making in-situ spectral
data truly usable---not merely archived.


% =========================================================================
\section*{Acknowledgements}
% =========================================================================

We acknowledge ISRO for the public release of Chandrayaan-3 LIBS data
through PDS4-compliant archives, and NIST for maintaining the Atomic
Spectra Database.  We thank the open-source scientific Python ecosystem
(NumPy, SciPy, Pandas, Matplotlib, Astropy) and the planetary science
community whose LIBS instrument and calibration work informs this study.


% =========================================================================
\section*{Author Contributions}
% =========================================================================

\textbf{L.~Anand:} Conceptualization, system architecture, algorithm
development, database design, API implementation, LTTB formulation,
data analysis, writing---original draft and editing.

\textbf{D.~Saeed:} Data processing implementation, PDS4 parsing,
LTTB validation, visualisation design, performance benchmarking,
writing---review and editing.


% =========================================================================
\section*{Data Availability}
% =========================================================================

Processed Chandrayaan-3 LIBS data products and LunarAtlas source code
will be made publicly available upon manuscript acceptance at a
Zenodo-hosted DOI.  Raw Chandrayaan-3 LIBS data are available from
ISRO's PDS4 archives at \url{https://pradan.issdc.gov.in/}.


% =========================================================================
\section*{Competing Interests}
% =========================================================================

The authors declare no competing interests.


% =========================================================================
%  References
% =========================================================================
\begin{thebibliography}{99}

\bibitem{Cremers2006}
D.~A. Cremers and L.~J. Radziemski,
\textit{Handbook of Laser-Induced Breakdown Spectroscopy}, 2nd~ed.,
John Wiley \& Sons, 2013.

\bibitem{Wiens2012}
R.~C. Wiens~\emph{et~al.},
``The ChemCam instrument suite on the Mars Science Laboratory (MSL)
rover: body unit and combined system tests,''
\textit{Space Sci. Rev.}, vol.~170, pp.~167--227, 2012.

\bibitem{Maurice2012}
S.~Maurice~\emph{et~al.},
``The ChemCam instrument suite on the Mars Science Laboratory (MSL)
rover: science objectives and mast unit description,''
\textit{Space Sci. Rev.}, vol.~170, pp.~95--166, 2012.

\bibitem{Gasnault2023}
O.~Gasnault~\emph{et~al.},
``ChemCam: zapping Mars for 10~years (and more),''
in \textit{54th Lunar and Planetary Science Conf.}, Abstract~\#2076,
2023.

\bibitem{Wiens2020}
R.~C. Wiens~\emph{et~al.},
``The SuperCam instrument suite on the NASA Mars~2020 rover: body
unit and combined system tests,''
\textit{Space Sci. Rev.}, vol.~217, article~89, 2020.

\bibitem{Xu2021}
W.~Xu~\emph{et~al.},
``The MarSCoDe instrument suite on the Mars rover of China's
Tianwen-1 mission,''
\textit{Space Sci. Rev.}, vol.~217, article~64, 2021.

\bibitem{Zhao2023}
Y.-Y. S. Zhao~\emph{et~al.},
``In situ analysis of surface composition and meteorology at the
Zhurong landing site on Mars,''
\textit{Natl. Sci. Rev.}, vol.~10, 2023.

\bibitem{ISROCh3}
Indian Space Research Organisation,
``Chandrayaan-3 Mission Updates,''
\url{https://www.isro.gov.in/Chandrayaan3.html},
accessed April 2026.

\bibitem{Laxmiprasad2020}
A.~S. Laxmiprasad~\emph{et~al.},
``Laser induced breakdown spectroscope on Chandrayaan-2 rover:
a miniaturized mid-UV to visible active spectrometer for lunar surface
chemistry studies,''
\textit{Curr. Sci.}, vol.~118, no.~4, pp.~573--581, 2020.

\bibitem{Saeidfirozeh2024}
H.~Saeidfirozeh~\emph{et~al.},
``Laser-induced breakdown spectroscopy in space applications:
Review and prospects,''
\textit{Trends Anal. Chem.}, vol.~181, article~117991, 2024.

\bibitem{Hughes2018}
S.~J. Hughes, A.~C. Raugh, and T.~L. Crichton,
``Planetary Data System Standards Reference,''
NASA Planetary Data System, version~1.14.0, 2018.

\bibitem{Steinarsson2013}
S.~Steinarsson,
``Downsampling time series for visual representation,''
M.Sc. thesis, University of Iceland, 2013.

\bibitem{Knight2000}
A.~K. Knight, N.~L. Scherbarth, D.~A. Cremers, and M.~J. Ferris,
``Characterization of laser-induced breakdown spectroscopy (LIBS)
for application to space exploration,''
\textit{Appl. Spectrosc.}, vol.~54, no.~3, pp.~331--340, 2000.

\bibitem{Lasue2012}
J.~Lasue~\emph{et~al.},
``Remote laser-induced breakdown spectroscopy (LIBS) for lunar
exploration,''
\textit{J. Geophys. Res. Planets}, vol.~117, E01002, 2012.

\end{thebibliography}

\end{document}